\documentclass[11pt]{article}
\usepackage[T1]{fontenc}
\usepackage{lmodern,amsmath,amssymb,geometry,enumitem,xcolor,hyperref}
\geometry{letterpaper,margin=.8in}
\hypersetup{colorlinks=true,urlcolor=blue!50!black,linkcolor=black}
\setlength{\parindent}{0pt}\setlength{\parskip}{5pt}
\setlist[itemize]{nosep,leftmargin=1.4em}
\newcommand{\R}{\mathbb R}\newcommand{\Z}{\mathbb Z}
\begin{document}
\begin{center}
{\Large An open pocket-detection question in rational billiards}\\[4pt]
Daniel Wang and Peyton Wang\\
\small Research note, 2 November 2026. Conjectural statements are identified below.
\end{center}

\section*{The question we would value help with}
Does every room in the class below have a nonsingular periodic billiard trajectory with nonzero response to the specified pocket slide? Is there an existing theorem that forces this response, or an obstruction that can occur in this planar room class?

This question is separate from the attached elementary construction of an orthogonal 22-wall room with a non-illuminating pair. That construction asserts existence and makes no minimality claim. The question here arose while attempting a reduction toward a lower bound; that reduction remains unfinished.

\section*{Room class and deformation}
Use the fixed lattice triangulation whose triangles are
\[
\operatorname{conv}\{c,c+u_j,c+u_{j+1}\},\qquad c\in(2\Z)^2,
\]
where $(u_0,\ldots,u_8)=((1,0),(1,1),(0,1),(-1,1),(-1,0),
(-1,-1),(0,-1),(1,-1),(1,0))$.
A room $P$ is a finite union of these triangles with connected interior and strictly simple polygonal boundary. It has no holes or internal mirror partitions. Require every integer point of $P$ outside $(2\Z)^2$ to be a physical boundary corner. We call this the \emph{raw condition}. The room contains two specified distinct interior points $S,T\in(2\Z)^2$.

Require the triangle with vertices $A=(2,0), B=(1,1), C=(1,0)$ to be a leaf pocket, attached through the open portal $BC$. Immediately above $B$ and below $C$, the boundary lies on the same line $x=1$. Translate the pocket by $t(0,1)$ for sufficiently small $t$. The adjoining vertical flanks change length along their fixed supporting line; all other mirror supporting lines stay fixed. The moving mirrors are
\[
 D_t:\ x+y=2+t,\qquad H_t:\ y=t.
\]
The two marked points stay fixed. Only the base room at $t=0$ is required to lie in the lattice triangulation.

\section*{Precisely what ``response'' means}
Fix a nonsingular periodic trajectory at $t=0$ that closes with its unfolding frame. Let $R_j$ be its successive linear reflection matrices and set $L_0=I$, $L_j=R_1\cdots R_j$. For its fixed reflection word define
\[
 Q(\gamma)=\sum_{j=1}^m L_{j-1}a_j,\qquad
 a_j=\begin{cases}(1,1)&\text{on }D,\\(0,2)&\text{on }H,\\(0,0)&\text{on a fixed mirror.}\end{cases}
\]
Up to an overall sign, $Q$ is the derivative of the unfolded period displacement under this slide. The question is whether some $\gamma$ has $Q(\gamma)\ne0$. We do not assert that the trajectory remains closed after sliding, or that nonzero $Q$ alone proves illumination of $S$ by $T$.

\newpage
\section*{What has been checked, and what has failed}
\begin{itemize}
\item A raw 25-wall room is blind in both direction classes $(1,1)$ and $(2,1)$, but has nonzero response in class $(3,2)$. It is obtained by flipping one distant triangular leaf in a previously studied room. Exact independent polygon checks cover every relevant strip in these directions.
\item The cancellation has a hand explanation. Before the flip, two pocket visits contribute $(1,-1)$ and $(-1,-1)$. The changed exterior return adds a visit contributing $(0,2)$, cancelling their total. An eight-reflection $(3,2)$ witness stays in the unchanged part of both rooms. This explanation has an independent hand audit.
\item A separate raw 37-wall room with 40 canonical triangles has $Q=0$ in every primitive direction class with $p^2+q^2<41$. Class $(5,4)$ has nonzero response. The geometry and complete fixed-strip coverage have independent exact verification. Directions are grouped under the reflection group $D_4$, with representatives $0\le q\le p$ and $\gcd(p,q)=1$.
\end{itemize}
Thus the previously proposed two-direction list, and the seven-class list through squared norm 25, fail. These examples do not exclude a larger fixed finite list or establish an unbounded first-detecting norm.

A parity shortcut is insufficient. For any closed core,
\[
 Q(\gamma)\equiv \#\{D\text{ hits}\}\,(1,1)\pmod2.
\]
The pocket seam determines a nonzero mod-2 cohomology class under the additional-diagonal-wall hypothesis. A nonzero class, however, need not be nonzero on a cylinder core of a general square-tiled surface. Abstract surface examples therefore prevent using connectedness of its double cover as a proof. They do not give a counterexample in the raw planar class.

\section*{The further gap concerning the actual marks}
Even a positive answer about periodic cores leaves an additional question. We seek a broken billiard path between the \emph{actual} interior marks $S,T$ through exactly one physical corner at its developed midpoint. Its two developed half-vectors must be equal and primitive, all other contacts must avoid physical corners, and the selected continuation must bound an embedded intrinsic angle of $\pi$ in the unfolded surface. Its transverse opening response under the slide must be nonzero.

A loaded periodic core does not automatically provide that marked event. Conditional affine-transfer formulas preserve an already supplied event, but require a nonzero geometric overlap contrast that is not known in general. Neither universal periodic-core detection nor this coupling to the actual marks is proved. No general lower bound of 22 walls follows from the present work.

\section*{Specific requests for mathematical input}
Is the planar raw condition strong enough to force a cylinder on which the pocket-slide class evaluates nontrivially? If so, what theorem or geometric invariant could supply this fact? Separately, what additional hypothesis could turn such a cylinder into the required event between the specified marked points?

\small
Exact polygons and contact certificates are available for the two examples above. Related published context is A.~Wolecki, \emph{Illumination in Rational Billiards}, arXiv:1905.09358v3 (2024), Theorem~1.3 and Remark~2.2, \url{https://arxiv.org/abs/1905.09358v3}. That work establishes the known unrestricted 22-edge upper bound; no implication for the deformation question above is claimed.
\end{document}
